\documentclass[10pt,letterpaper]{article}

\usepackage[letterpaper,margin=0.82in]{geometry}
\usepackage{amsmath,amssymb}
\usepackage{array,booktabs,tabularx}
\usepackage[hidelinks]{hyperref}

\hypersetup{
  pdftitle={WKVM: Treating Model State as a First-Class Serving Object},
  pdfauthor={xiaol}
}

\setlength{\emergencystretch}{2em}
\setlength{\parskip}{0.15em}
\renewcommand{\arraystretch}{1.08}
\newcolumntype{Y}{>{\raggedright\arraybackslash}X}
\newcommand{\code}[1]{\texttt{#1}}
\newcommand{\wkvm}{\textsc{WKVM}}

\title{\textbf{WKVM: Treating Model State as a First-Class Serving Object}\\
\large State-Native Inference for Recurrent and Hybrid Language Models}
\author{xiaol}
\date{Technical Report --- August 2026}

\begin{document}

\maketitle

\begin{abstract}
Contemporary language-model servers are organized around a key--value (KV)
cache whose size grows with the retained token history.  Recurrent and linear
attention models expose different serving physics: each request carries a
fixed-size execution state, independent of context length.  We present
\wkvm, a prototype that makes this state the primary allocation and lifecycle
object.  The implemented native RWKV-7 path combines typed state families,
all-or-nothing slot admission, a continuous-batching scheduler over
computed-token gaps, dense GPU state banks, and named versioned snapshots that
can be restored, forked, persisted, and mutated with lineage.  A separate Gemma
guest path retains a bounded set of actual KV spans selected by value-vector
routing; it is approximate and is not an RWKV recurrence.

We report the evidence with its limitations.  On one RTX 4090, a single-run
RWKV-7 1.5B steady-state experiment reached 8,077 tokens/s at batch 256 with a
12.19 MiB state per request.  A separate 191M state-store demonstration created
2,000 snapshots at 2.29 MiB each, measured 8.2 ms median warm restore and 8.6 ms
median immediately reloaded disk-backed restore, and reproduced 16 of 16
interrupted greedy continuations exactly.  For bounded Gemma serving, synthetic
recall averaged 0.90 versus 0.95 for full KV over 8K--32K contexts.  In three
clean A800 runs at batch 64 and 16K input tokens, \wkvm{} used a reported
28.4 GiB whole-GPU peak and delivered 4.47$\times$ the comparable vLLM decode
rate, but only 0.79$\times$ its end-to-end output throughput because prefill
dominated.  These results support fixed resident state and explicit state
lifecycle management; they do not support quality equivalence or a universal
speedup claim.
\end{abstract}

\section{Introduction}

High-throughput language-model serving has been shaped by autoregressive
Transformers.  Systems such as Orca, vLLM, and SGLang schedule token-level work
continuously and manage large per-request KV histories through iteration-level
batching, paging, and prefix reuse~\cite{orca,vllm,sglang}.  This design is a
good match for full attention, where the history is both large and naturally
indexed by a token prefix.  It is not the only useful model-state regime.
RWKV, Mamba, Gated DeltaNet, and related recurrent or linear-attention models
compress history into a fixed-size state updated at every step
~\cite{rwkv,rwkv7,mamba,gateddeltanet}.

The distinction is more than a memory optimization.  A compact recurrent state
can be treated as a session object: it can be parked, restored, copied for a
branch, exported, or deliberately transformed.  A token-prefix cache usually
assumes that the stored object is derived from a prefix, $S=f(x_{1:t})$.
Controlled mutation instead produces $S' = g(S)$, whose identity and provenance
cannot be expressed by the token prefix alone.  Supporting that operation does
not make prefix caching impossible; it requires a separate handle, version, and
lineage layer.

\wkvm---``WKV plus KVM''---explores this state-native design.  Its current
implementation contains two related but distinct systems:

\begin{enumerate}
  \item a native pure-RWKV-7 execution path whose primary allocation unit is a
  fixed slot for each state family; and
  \item a Gemma-4 guest path that bounds memory by retaining and routing actual
  KV spans.  This second path is a training-free approximation, not a recurrent
  matrix-state implementation.
\end{enumerate}

This separation matters.  Earlier project design notes discuss broader GDN,
Mamba, and hybrid recurrent banks, but the present native loader rejects hybrid
RWKV checkpoints, and those broader backends are future work.  The report
therefore makes four narrower contributions:

\begin{itemize}
  \item \textbf{A slot-native substrate.}  Models declare typed fixed-footprint
  state families.  Admission allocates one slot in every family atomically, and
  the scheduler advances a single computed-token invariant for prefill, decode,
  and resume.
  \item \textbf{An explicit state lifecycle.}  Named, versioned records bind
  state bytes to token/accounting metadata and lineage.  Pinned-host and
  safetensors-backed tiers support snapshot, restore, fork, persistence, and
  registered mutations in the prototype.
  \item \textbf{A bounded transformer case study.}  Gemma routed-span mode
  combines an exact recent ring with content-routed retained spans and an
  application-level parent-token contract for safe multi-turn reuse.
  \item \textbf{A claim-audited evaluation.}  We retain negative results,
  distinguish repeated from single-run evidence, and separate exact RWKV
  semantics from approximate Gemma comparisons.
\end{itemize}

The central conclusion is correspondingly modest: treating compact model state
as a first-class object simplifies capacity accounting and enables lifecycle
operations that are awkward under prefix identity alone.  It does not remove
prefill cost, state bandwidth, model-quality limits, or the need for mature
kernels.

\section{Background and Scope}

\subsection{KV histories and recurrent state}

For a Transformer with retained length $L$, the KV footprint is approximately
\begin{equation}
  M_{\mathrm{KV}}(L)
  = 2L b \sum_{\ell \in \mathcal{A}} h_{\ell}^{\mathrm{kv}} d_{\ell},
  \label{eq:kv}
\end{equation}
where $b$ is bytes per element and $\mathcal{A}$ is the set of attention
layers.  Sliding-window and KV-sharing layers change the coefficient, but full
attention remains length-dependent.  PagedAttention reduces fragmentation and
enables sharing without changing this semantic scaling~\cite{vllm}.

A recurrent layer instead carries an update state
\begin{equation}
  S_t = F(S_{t-1}, x_t), \qquad y_t = G(S_t, x_t),
  \label{eq:recurrent}
\end{equation}
whose storage is fixed for a given model.  In the implemented RWKV-7 path, the
state contains an FP32 matrix-valued WKV component and two model-dtype
token-shift vectors per layer.  The serving footprint is therefore
\begin{equation}
  M_{\mathrm{state}} = \sum_{j=1}^{J} m_j,
  \label{eq:state}
\end{equation}
where $m_j$ is the fixed footprint of state family $j$; it does not depend on
$L$.

Constant context scaling does not imply negligible cost.  Linear-attention
states can be large enough that reading and writing them every token becomes a
bandwidth bottleneck.  KVBuffer directly targets this issue by buffering recent
K/V and amortizing recurrent-state updates~\cite{kvbuffer}.  \wkvm's slot and
lifecycle abstractions are complementary to such kernel-level update policies.

\subsection{Identity, reuse, and lifecycle}

PagedAttention primarily addresses physical allocation; RadixAttention adds
prefix-indexed reuse~\cite{vllm,sglang}.  Marconi checkpoints recurrent-layer
state for hybrid-model prefix caching~\cite{marconi}.  Pensieve and other
stateful-session systems retain and tier multi-turn Transformer KV
histories~\cite{pensieve}.  These systems establish that state reuse and session
persistence are not new by themselves.

\wkvm{} focuses on a narrower identity boundary.  A prefix-derived checkpoint is
immutable evidence of $f(x_{1:t})$.  A first-class recurrent-state handle may
also represent a controlled transform, merge, or decay whose parent and rule
must be recorded independently of the visible token sequence.  The prototype
does not establish that arbitrary mutations preserve language-model quality;
it supplies the mechanism and provenance needed to study them.

\subsection{Goals and non-goals}

The implemented goals are exact slot-count admission, continuous arrivals,
fixed-footprint recurrent state ownership, explicit hibernate/resume, and a
bounded-memory Gemma compatibility path.  The current system is not a general
RWKV/GDN/Mamba server, a drop-in replacement for the complete OpenAI API, a
distributed migration service, or a quality-equivalent replacement for full
KV Gemma.  CUDA-graph results in this report come from a benchmark wrapper that
captures the model forward for fixed batch buckets; the core RWKV runner still
performs scheduler work and state gather/scatter eagerly.

\section{State-Native Architecture}

\begin{figure}[t]
\centering
\fbox{\begin{minipage}{0.94\linewidth}
\centering
\small
token-ID frontend $\longrightarrow$ gap scheduler
$\longrightarrow$ state arena $\longrightarrow$ model runner\\[4pt]
$c_r\rightarrow n_r$ \hspace{20mm} family slot IDs
\hspace{17mm} prefill / decode\\[7pt]
dense GPU state banks $\longleftrightarrow$ pinned-host versioned records
$\longleftrightarrow$ safetensors + JSON index
\end{minipage}}
\caption{Implemented native RWKV-7 data path.  The arena owns integer slot
identities, the GPU bank owns tensors, and the optional store owns immutable
snapshot records.  Transfers are synchronous in the current prototype.}
\label{fig:architecture}
\end{figure}

\subsection{Typed state families and exact admission}

A model exposes a \code{ModelStateSpec} containing one or more
\code{StateFamilySpec} records.  Each family declares a name, a byte footprint
per request, and the layers that use it.  The RWKV-7 implementation declares a
\code{wkv} family and a \code{shift} family.  The arena preallocates integer
slot IDs $1,\ldots,N$ in every family; slot zero is reserved as a padded write
target for static batch shapes.

For request $r$, allocation returns
\begin{equation}
  A(r) = \{s_{r,1},\ldots,s_{r,J}\},
\end{equation}
with one slot from every family.  Admission is possible exactly when
\begin{equation}
  \min_j |\mathcal{F}_j| > 0,
  \label{eq:admission}
\end{equation}
where $\mathcal{F}_j$ is the free set for family $j$.  This removes
token-block fragmentation and watermark reasoning from recurrent-state
admission.  It trades that simplicity for static provisioning: unused bytes
inside a family slot cannot be lent to another footprint class.

The tensor owner is separate from the allocator.  \code{RWKV7StateBank}
materializes layer-major arrays so that a layer's active batch is gathered with
one indexed selection and committed with one indexed copy.  Python request
objects carry only slot IDs, not state tensors.

\subsection{One computed-token scheduling invariant}

Each request stores prompt tokens, sampled output tokens, and a single counter
$c_r$ for tokens whose state contribution has been computed.  Its current
target is
\begin{equation}
  n_r = |x_r| + |y_r|, \qquad g_r = n_r-c_r.
  \label{eq:gap}
\end{equation}
A request with $g_r=1$ is in steady-state decode.  A request with a larger gap
has prefill or continuation work.  A restored request simply begins with
$c_r>0$ and a nonempty gap; it does not require a separate scheduler phase.

At every step, the scheduler first advances running requests under a global
token budget, then admits waiting requests while slots and running capacity
remain.  A per-request cap prevents a long prompt from consuming the entire
step.  Completed requests release their slots during the same update, allowing
new arrivals to enter on the next step.  A completion-prefill lane is also
implemented for workload profiles that favor finishing a bounded cohort of
long prompts while preserving decode priority.

\subsection{RWKV execution and graph buckets}

The native loader accepts FLA-format pure RWKV-7 checkpoints.  Prefill seeds an
FLA cache from arena slots, executes chunks, and writes the resulting recurrent
and shift states back after each chunk.  Decode gathers a changing set of slot
rows, performs one recurrent step for the batch, and scatters the state back.
The reference FLA kernels provide chunked prefill and fused recurrent decode;
\wkvm{} owns scheduling and persistent state placement~\cite{fla}.

The M2 benchmark adds one CUDA graph for the model forward at each measured
batch size.  Static graph input and output buffers are replayed when the batch
shape is unchanged.  State gather/scatter, sampling, and scheduler Python
remain eager, so these results should not be described as an integrated
whole-engine graph implementation.

\subsection{Versioned durable-state records}

An optional \code{StateStore} snapshots a live GPU slot into pinned host
memory.  A record can be summarized as
\begin{equation}
  H = (\text{name},\text{version},\phi,c,\tau,
       \text{parent},\text{rule},\theta),
  \label{eq:handle}
\end{equation}
where $\phi$ is the current state-layout fingerprint, $c$ is the computed-token
count, $\tau$ is the known token list, and $\theta$ stores rule parameters.
Handles have the form \code{name@version}; versions are append-only and records
are immutable after creation.

The prototype exposes the following operations:

\begin{itemize}
  \item \textbf{snapshot/hibernate}: copy state from the GPU bank to pinned
  host memory, optionally releasing the live slot;
  \item \textbf{restore}: allocate a slot, transfer the saved tensors to it,
  and submit a request whose $c_r$ starts at the saved count;
  \item \textbf{fork}: create a new metadata record sharing immutable warm
  tensors until a transformed child is materialized;
  \item \textbf{persist}: write contiguous tensors to one safetensors file and
  update an atomically replaced JSON index; and
  \item \textbf{mutate}: clone the parent tensors, apply a registered rule such
  as decay or merge, and record parent/rule/parameters in a new version.
\end{itemize}

The present fingerprint hashes shapes and dtype, not model weights or a full
checkpoint manifest.  The store is synchronous and assumes a single-threaded
engine.  Warm-only snapshots do not survive process death; only explicitly
persisted records can be rebuilt from the index.  These constraints are
important boundaries for any production interpretation.

\section{Bounded Gemma Guest Mode}

\subsection{Actual retained-KV design}

The Gemma guest is a separate execution engine.  It constructs the text model
from checkpoint tensors and includes custom scheduling, token-pool, and Triton
decode paths, but it still uses Transformer configuration and tokenizer
utilities at the server boundary.  More importantly, its bounded memory is
made of retained \emph{actual keys and values}; it is not the RWKV-7 matrix
state used by the native path.

For each supported full-attention layer, a request keeps an exact sink, an
exact recent ring, a pending eviction region, and a fixed number of routed
span slots.  Evicted text is split at punctuation with fallbacks, a leader
layer selects a span using its most novel value-vector feature, and online
cosine routing assigns it to a persistent slot.  Within a slot, farthest-point
diversity retention prunes spans under a token budget.  The commonly evaluated
configuration uses sink 16, ring 1,024, pending 512, 64 routed slots, and a
144-token limit per slot.  Accounting in the evidence bundle bounds retained
material at approximately 10,831 tokens per row for that configuration.

This is a fixed-capacity compression policy.  Recent dependencies inside the
ring are exact; information outside it can be discarded, coalesced, or routed
to a competing slot.  It therefore belongs with bounded/streaming attention
methods such as StreamingLLM and KV-cache compression, not with exact
Transformer serving~\cite{streamingllm,snapkv}.

\subsection{Parent-token session contract}

The Open WebUI integration retains a parked engine session after a turn and
reuses it only when a provider-specific parent-token contract validates.  The
binding includes the model, user, chat, current and parent message IDs, exact
visible parent history, and a digest of raw retained token IDs.  This matters
because decoding generated tokens to text and re-tokenizing is not always an
identity operation.  Edits, branches, stale parents, changed filters, or
missing metadata cause a safe fresh prefill rather than silent reuse.

The HTTP surface is deliberately narrow: text-only greedy chat, one choice,
and no tools, images, log probabilities, or arbitrary stop sequences.  The
regular server binds to loopback and has no built-in API-key enforcement.

\section{Evaluation}

\subsection{Questions and evidence policy}

We ask four questions:

\begin{enumerate}
  \item[\textbf{Q1.}] Does the native RWKV path preserve a fixed per-session footprint while
  supporting large decode batches?
  \item[\textbf{Q2.}] Do lifecycle operations restore exact continuations at useful latency?
  \item[\textbf{Q3.}] What recall is lost by bounded routed-span Gemma memory?
  \item[\textbf{Q4.}] Do bounded resident sessions translate into end-to-end serving wins?
\end{enumerate}

The experiments were produced during rapid development between July 3 and
July 23, 2026.  We distinguish repeated, clean campaigns from descriptive
single-run checkpoints.  Cross-engine Gemma rows also compare different
semantics: \wkvm{} uses \code{routed\_span\_approximate}, while vLLM and SGLang
use full KV.  Such ratios characterize these concrete systems and workloads;
they are not quality-normalized engine rankings.

\subsection{Q1: Native RWKV decode}

Table~\ref{tab:rwkv} transcribes the M2 steady-state benchmark on one RTX 4090.
Weights are BF16 and WKV state is FP32.  Every timed token passes through the
engine step, but each cell is one 64-step timing after eight warmup steps; raw
per-repeat data and confidence intervals are unavailable.

\begin{table}[t]
\centering
\caption{Selected native RWKV-7 steady-state results on one RTX 4090.  Graphs
capture the fixed-shape model forward, not the complete engine step.}
\label{tab:rwkv}
\small
\begin{tabular}{@{}lrrrrr@{}}
\toprule
Model & State/slot & Batch & Eager tok/s & Graphed tok/s & Peak VRAM \\
\midrule
1.5B & 12.19 MiB & 1   & 67    & 195   & 2.96 GiB \\
1.5B & 12.19 MiB & 64  & 3,592 & 5,274 & 5.92 GiB \\
1.5B & 12.19 MiB & 256 & 7,913 & 8,077 & 19.02 GiB \\
191M & 2.29 MiB  & 1   & 132   & 680   & 0.45 GiB \\
191M & 2.29 MiB  & 256 & 26,707 & 38,856 & 2.94 GiB \\
\bottomrule
\end{tabular}
\end{table}

The result demonstrates the intended footprint and batchability.  Graph
benefit is largest at batch one (2.91$\times$ for 1.5B) and falls to
1.02$\times$ at batch 256, where model work dominates.  The comparison does
not establish kernel leadership: the runner uses generic FLA kernels and eager
state copies, and the experiment includes no matched incumbent serving path
for the same checkpoint.

\subsection{Q2: State lifecycle}

The M3 demonstration uses RWKV-7 191M on the same GPU.  It creates 2,000
sessions through continuous batches of 64, snapshots each to pinned host
memory, and persists 500 records.  Table~\ref{tab:store} reports descriptive
single-run results.

\begin{table}[t]
\centering
\caption{State-store demonstration.  ``Disk-backed reload'' was measured
immediately after persistence without flushing the operating-system page
cache; it is not a cold-device benchmark.}
\label{tab:store}
\small
\begin{tabularx}{\linewidth}{@{}YrrY@{}}
\toprule
Operation & Samples & Reported result & Interpretation \\
\midrule
Pinned-host restore to next token & 224 & median 8.2 ms & synchronous H2D restore \\
Disk-backed immediate reload & 76 & median 8.6 ms & page cache may dominate \\
Interrupted/resumed exactness & 16 & 16/16 identical & greedy, batch-one twins \\
Fork creation & 64 & 544.8 ms total & 8.5 ms/fork; immutable bytes shared \\
Restart recovery & 1 session & 24-token match & explicitly persisted record \\
Decay mutation & 1 lineage & continuation changed & mechanism demo, not quality test \\
\bottomrule
\end{tabularx}
\end{table}

The original artifact also reports tail percentiles, but the benchmark's
index-based percentile calculation is not a statistically sound p99 estimate
for these sample counts, so we exclude those claims.  The exactness check is
more informative: it shows that, for the tested greedy path and state layout,
snapshot transfer does not change continuation tokens.  It does not test
cross-checkpoint rejection because the current fingerprint omits weight
identity.

\subsection{Q3: Routed-span recall}

The routed-span quality grid evaluates Gemma-4-E4B-it~\cite{gemma4model} at 8,192, 16,384, and
32,768 tokens on three synthetic tasks: one needle, one of eight keyed facts,
and aggregation of six items.  Generation is greedy and substring-scored.
Needle and multi-key cells average three seeds; aggregation uses one seed.
Table~\ref{tab:quality} reports means over 15 task/depth cells per context.

\begin{table}[t]
\centering
\caption{Synthetic recall means for Gemma bounded-memory variants.}
\label{tab:quality}
\small
\begin{tabular}{@{}lrrrr@{}}
\toprule
Mode & 8K & 16K & 32K & Overall \\
\midrule
Full KV & 0.97 & 0.95 & 0.94 & 0.95 \\
Sink + recent ring & 0.18 & 0.00 & 0.00 & 0.06 \\
Temporal segment bank & 0.51 & 0.39 & 0.44 & 0.45 \\
Value-routed slots & 0.77 & 0.71 & 0.90 & 0.79 \\
Span-atomic value routing & 0.89 & 0.91 & 0.90 & 0.90 \\
\bottomrule
\end{tabular}
\end{table}

Span-atomic routing improves the interference-heavy multi-key task from 0.51
to 0.89 overall by keeping fact spans together.  It slightly reduces the
scattered aggregation score (0.86 to 0.81), and its 32K retained memory is
reported at roughly 88 MiB per slot versus 36 MiB for the ring-only mode.  The
grid is evidence for this synthetic retrieval mechanism, not general language
quality.  Historical natural-document NLL increases outside the exact window,
and native natural-document NLL has not been rerun.

\subsection{Q4: Serving comparisons}

The strongest repeated comparison is a clean three-run campaign on one
NVIDIA A800-SXM4-80GB.  All engines receive 64 uniform 16,384-token synthetic
prompts and generate 32 fixed-length greedy tokens.  Results are medians, with
min--max ranges in the source artifact.

\begin{table}[t]
\centering
\caption{Repeated A800 B64, 16K-input, 32-output campaign.  Memory is the
reported whole-GPU peak and is affected by each incumbent's configured pool
fraction; it is not minimum required memory.}
\label{tab:a800}
\small
\begin{tabular}{@{}lrrrr@{}}
\toprule
Engine & Semantics & E2E tok/s & Decode tok/s & Peak used \\
\midrule
\wkvm{} & routed-span approx. & 21.886 & 121.612 & 28.438 GiB \\
vLLM 0.25.1 & full KV & 27.608 & 27.227 & 74.743 GiB \\
SGLang 0.5.15 & full KV & 15.587 & excluded & 69.930 GiB \\
\bottomrule
\end{tabular}
\end{table}

Against vLLM, \wkvm{} reaches 4.467$\times$ the comparable decode rate but only
0.793$\times$ end-to-end throughput.  Its cohort prefill wall is 84.8 s versus
73.6 s for vLLM, and prefill dominates a workload with only 32 output tokens.
Against SGLang, \wkvm{} is 1.404$\times$ on end-to-end output throughput; SGLang's
decode estimate is excluded because it is based on separate-run subtraction.
This campaign is the clearest counterexample to a universal speedup claim.

An RTX 4090 B16 audit gives a complementary context-scaling result.  At 16K
input and 128 output tokens, three \wkvm{} runs average 64.769 E2E tokens/s,
while one controlled vLLM run reaches 64.713---an E2E tie.  Comparable decode is
256.467 versus 67.091 tokens/s, and the mean whole-GPU engine delta is 16.871
versus 18.405 GiB.  When context doubles to 32K, \wkvm's single-run decode
remains 253.746 tokens/s and memory increases by only 0.055 GiB, while prefill
roughly doubles.  The archived 32K incumbent rows use a different idle GPU
baseline, so their E2E ratios remain exploratory.

Finally, an application-level Open WebUI checkpoint drives 32 persisted
conversations through eight synchronized turns.  It completes 256 requests at
345.097 output tokens/s with 224/224 eligible continuations reused, versus
160.322 for the tested vLLM profile and 65.055 for the tested SGLang profile.
This 2.153$\times$/5.305$\times$ result is a single cross-run checkpoint with
engine-dependent generated histories, approximate-versus-full semantics, and
external raw request artifacts.  It demonstrates the parent-token lifecycle
path; it is not a repeated publication envelope.

\section{Discussion}

\subsection{What the abstraction buys}

For a fixed-footprint recurrent model, slot-native scheduling makes admission
an exact count rather than a token-block packing problem.  Dense state banks
also make request composition an index-selection problem, which is compatible
with stable batch buckets.  Most importantly, a named state handle separates
session identity from a token prefix.  Forking, explicit persistence, and
provenance-bearing mutation can therefore be API operations instead of hidden
cache behavior.

This design complements rather than replaces prior work.  Marconi decides when
prefix-derived recurrent checkpoints are worth caching; KVBuffer decides when
to flush recurrent state for better IO efficiency; Pensieve tiers attention KV
for multi-turn sessions~\cite{marconi,kvbuffer,pensieve}.  A future mature
runtime could combine all three ideas: typed slot allocation, checkpoint-aware
prefix reuse, and buffered state updates.

\subsection{Where fixed state does not help}

Fixed resident memory does not make prompt computation constant-time.  Both
the A800 and 4090 campaigns expose prefill as the principal end-to-end
bottleneck.  Recurrent state itself can also be bandwidth-heavy, as the
diminishing graph benefit at high batch and KVBuffer's analysis emphasize.
At short context, low concurrency, or short sessions, a mature full-KV engine
may simply be faster.

The Gemma guest makes an additional trade: it bounds memory by discarding or
coalescing history.  Its recall score is workload-dependent, and exact full-KV
systems retain a semantic advantage.  Performance comparisons must therefore
publish the semantics column, not only a throughput ratio.

\subsection{Mutable state as a research capability}

The built-in decay and merge rules demonstrate lineage, not validated memory
editing.  A transformed recurrent state may be useful for consolidation,
forgetting, branching agents, or post-training experiments, but it may also
produce arbitrary behavior.  Necessary next steps include checkpoint-bound
fingerprints, rule-specific invariants, reference recomputation where
available, downstream quality evaluation, and access control for state
capabilities.

\section{Related Work}

\paragraph{Transformer serving.}
Orca introduced iteration-level scheduling for generative models~\cite{orca}.
vLLM's PagedAttention reduces KV fragmentation and supports sharing, while
SGLang's RadixAttention indexes reusable KV by token-prefix paths
~\cite{vllm,sglang}.  Llumnix studies live migration and scheduling across LLM
instances~\cite{llumnix}.  Pensieve retains and tiers multi-turn KV state
~\cite{pensieve}.  These systems motivate \wkvm's lifecycle vocabulary but
operate primarily on attention KV rather than opaque mutable recurrent state.

\paragraph{Recurrent and hybrid serving.}
RWKV, Mamba, and Gated DeltaNet establish fixed-state sequence models and
parallel/recurrent computation forms~\cite{rwkv,rwkv7,mamba,gateddeltanet}.
Marconi caches prefix-derived recurrent checkpoints in hybrid LLMs
~\cite{marconi}; KVBuffer changes the IO schedule for linear state updates
~\cite{kvbuffer}; STree studies speculative tree verification for hybrid state
space models~\cite{stree}.  \wkvm{} differs in emphasis: the recurrent state is an
explicit versioned lifecycle object and the allocator's primary unit.  It does
not claim the first recurrent-state server or cache.

\paragraph{Bounded and growing memory.}
StreamingLLM preserves attention sinks and a recent window for constant-memory
streaming, while SnapKV compresses prompt KV based on attention-derived
importance~\cite{streamingllm,snapkv}.  Memory Caching stores recurrent hidden
checkpoints and aggregates selected memories to extend RNN capacity
~\cite{memorycaching}.  Gemma routed-span belongs to this approximate-memory
family: it retains selected actual KV spans under a fixed budget.  The native
RWKV state-store contribution is logically separate.

\section{Limitations and Threats to Validity}

\begin{itemize}
  \item \textbf{Prototype scope.}  Native execution currently supports pure
  RWKV-7 FLA-format checkpoints.  General Mamba, GDN, and hybrid state families
  are design targets, not implemented backends.
  \item \textbf{Rapid source evolution.}  Results span multiple commits over
  three weeks.  The artifact appendix freezes the report and file hashes, but a
  single current binary did not produce every table.
  \item \textbf{Limited native evaluation.}  RWKV throughput and lifecycle
  experiments are single-run on one RTX 4090 and small 191M/1.5B checkpoints.
  They lack uncertainty estimates, energy measurements, and matched engines on
  identical RWKV weights.  The historical M3 log also warns that the FLA path
  still needs an official-RWKV implementation cross-check.
  \item \textbf{Store methodology.}  The disk-backed latency did not evict the
  OS page cache; reported p99 calculations are excluded; the fingerprint binds
  layout rather than checkpoint weights; copies synchronize; and the store is
  single-threaded.
  \item \textbf{Quality coverage.}  Routed-span quality uses small synthetic,
  substring-scored grids.  Native natural-document NLL and standard long-context
  benchmarks are missing.  Approximate and full-KV Gemma are not semantically
  equivalent.
  \item \textbf{Comparison statistics.}  Only the A800 cell has three clean
  runs for every engine.  The 4090 controlled vLLM point, 32K cells, and Open
  WebUI checkpoint need interleaved repeats, randomized order, confidence
  intervals, thermal/clock controls, and committed request-level artifacts.
  \item \textbf{Memory interpretation.}  Whole-GPU peaks include configured
  incumbent pools and any process on the selected device.  They demonstrate
  observed residency under those launch profiles, not minimum memory required
  by each engine.
\end{itemize}

\section{Artifact Availability and Reproducibility}

Source, benchmark drivers, result summaries, and provenance bundles are
available in the project repository~\cite{wkvmrepo}.  The dependency-free
allocator, scheduler, mixed-batch metadata, benchmark contract, and packaging
subset passes 60/60 tests in the report-preparation environment.  The full
794-test discovery run cannot be treated as a pass in that environment because
124 Torch-dependent tests error at import and 153 tests are skipped; historical
GPU artifacts were not rerun for this report.

Table~\ref{tab:artifacts} lists the exact evidence promoted into the paper.  A
companion claim ledger in \code{paper/CLAIMS.md} records excluded claims and
additional caveats.

\begin{table}[t]
\centering
\caption{Promoted artifact ledger.  SHA-256 values identify the report files,
not every external raw measurement referenced by those reports.}
\label{tab:artifacts}
\scriptsize
\begin{tabularx}{\linewidth}{@{}lYY@{}}
\toprule
Evidence & Repository path & Report SHA-256 \\
\midrule
RWKV M2 & \path{experiments/results/m2_engine_bench.md} &
\code{899d26677dadcf27...e508df7} \\
State lifecycle M3 & \path{experiments/results/m3_results.md} &
\code{15d336e000c7c4a...b951f9} \\
Routed-span quality & \path{experiments/results/quality_grid.md} &
\code{a5826b978281ecb...4f5c2} \\
4090 B16 audit & \path{experiments/results/gemma_b16_evidence_audit_20260713.md} &
\code{b5eef0bf411ca5b...0fae46} \\
A800 repeats & \path{experiments/results/gemma_a800_reliable_20260716/report.md} &
\code{83371ce4686212c...386776} \\
Open WebUI checkpoint & \path{experiments/results/open_webui_parent_token_b32_t8_20260723.md} &
\code{aea5cce4adcba99...a908} \\
\bottomrule
\end{tabularx}
\end{table}

\section{Conclusion}

\wkvm{} demonstrates that a serving runtime can be organized around fixed-size
model state rather than treating every architecture as a variant of a growing
KV history.  The implemented substrate provides exact family-slot admission, a
unified computed-token scheduler, dense recurrent state ownership, and
versioned lifecycle handles.  Its separate Gemma case study shows why bounded
resident state can improve long-session decode and memory behavior, and also
why prefill and approximation quality remain load-bearing constraints.

The current evidence is encouraging but not a universal performance claim.
The most useful next milestone is not another isolated speedup: it is a
checkpoint-bound, asynchronous state store plus interleaved repeated evaluation
on a mature native recurrent or hybrid model, including standard quality and
end-to-end workload curves.

\appendix

\section{Reproduction Entry Points}

The following commands correspond to the promoted artifact families.  GPU
commands require their documented model paths, accepted model licenses, and
environment-specific dependencies.

\begin{verbatim}
# Dependency-free current-checkout validation
python -m unittest -v \
  tests.test_arena tests.test_scheduler tests.test_mixed_batch \
  tests.test_benchmark_contract tests.test_packaging

# Historical native RWKV measurements
python experiments/m2_engine_bench.py --model /path/to/rwkv7 --graph
python experiments/m3_demos.py --model /path/to/rwkv7

# Frozen benchmark runners (see each script's --help and report contract)
bash scripts/run_wkvm_phase3_4090.sh
bash scripts/run_wkvm_10x_http_a800.sh
\end{verbatim}

The A800 repeated report records benchmark commit
\code{234cc04867a93f1352d2a1c220c216b698f46560}, model-manifest SHA-256,
source/worktree identity, GPU UUID, driver and package versions, launch
commands, request counts, output fingerprints, and min/median/max aggregates.
The sealed 4090 B16 bundle likewise contains copied JSON artifacts, source
manifests, commands, and integrity hashes.  Open WebUI raw request artifacts
remain external and are therefore not promoted as a repeatable envelope.

\section{Submission-Scope Statement}

This technical report intentionally excludes the following statements:
``\wkvm{} is 10$\times$ faster than vLLM/SGLang,'' full-quality equivalence for
routed-span Gemma, first stateful LLM serving, first recurrent prefix caching,
first linear-attention serving, production-ready live migration, and general
Mamba/GDN/hybrid support.  Any later revision that adds one of these claims
should add a new controlled artifact and update the claim ledger before
submission.

\sloppy
\bibliographystyle{unsrt}
\bibliography{references}

\end{document}
